\documentclass[10pt,a4paper]{amsart}
\usepackage[margin=19mm]{geometry}
\usepackage{amsmath,amssymb,mathtools}
\usepackage[scr=rsfs,cal=euler]{mathalfa}
\usepackage{xfrac}
\usepackage[unicode,hidelinks,bookmarks=false]{hyperref}
\newcommand{\ie}{i.e.\ }
\newcommand{\cf}{cf.\ }
\newcommand{\resp}{respectively\ }
\newcommand{\Subsection}{\subsection}
\providecommand{\nobreakdash}{\nobreak\hbox{-}}
\setlength{\emergencystretch}{2em}
\multlinegap0pt
\usepackage{amssymb}
\usepackage{stackengine}
\usepackage{xcolor}
\newtheorem{lemma}{Lemma}
\newcommand\Gal{{\mathrm {Gal}}}
\newcommand{\Q}{\mathbb{Q}}
\newcommand{\Z}{\mathbb{Z}}
\title{Euclidean Algorithms for Ideal Classes in Biquadratic fields: A Genus-Theoretic Perspective}
\author{Sunil Kumar Pasupulati}
\date{Source: arXiv:2512.00337v2 (CC BY 4.0)}
\begin{document}
\maketitle

\noindent\emph{Source and licence.} Euclidean Algorithms for Ideal Classes in Biquadratic fields: A Genus-Theoretic Perspective. Version arXiv:2512.00337v2 (CC BY 4.0), distributed by arXiv under the Creative Commons Attribution 4.0 International licence, http://creativecommons.org/licenses/by/4.0/, as recorded in the arXiv licence metadata for this exact version and shown on the abstract page https://arxiv.org/abs/2512.00337v2. Full licence notice: \texttt{licenses/arxiv-2512.00337-CC-BY-4.0.txt}.\par\smallskip

\noindent\textit{Editorial scope.} Counted unit: the article's lemma existence\_of\_lift3 (source0.tex lines 621--636). The article's complete original proof of it is delivered with it (source0.tex lines 638--656, one \texttt{proof} environment of 19 lines). No mathematics is written, repaired or re-derived: the statement and the proof are the article's own lines, in the article's own notation, and no retained line is substituted or re-flowed.  No further source-local context is needed: every cross-reference of every retained line resolves inside the two delivered spans.  Nothing was omitted from any retained span, so the package carries no omission record.  No retained line contains a citation, so the wrapper carries no bibliography and no bibliography entry.  No retained line contains a figure environment, an image-inclusion command or drawing code, so no image is delivered and the package declares no asset dependency.  Editorial formatting only, in addition: an AMS \texttt{amsart} preamble that restates the article's own declarations and macros used by the retained lines, namely the environments \texttt{lemma} on separate counters, together with the standard packages those lines need.  The retained environments print in this document as Lemma 1 in order of appearance, whereas the article prints the same items with the numbering of its own full document.  The complete native source of the article as distributed in the arXiv e-print for this exact version is bundled unchanged as \texttt{sources/190043/source0.tex}.
\begin{lemma}\label{existence_of_lift3}
Let $p_i$ be distinct primes with $p_3,p_4\equiv3\pmod4$, and let
\[
K=\Q(\sqrt{p_1p_2},\sqrt{p_3p_4})
\]
be a biquadratic field of class number two.  
If $\Gal(H(K)/K)=\langle\sigma\rangle$, then there exists a lift
$\hat{\sigma}\in\Gal(\Q(\zeta_f)/K)$ of $\sigma$ such that
\[
\hat{\sigma}\notin 
\bigcup_{\ell} \Gal(\Q(\zeta_f)/\Q(\zeta_\ell))
\;\cup\;
\Gal(\Q(\zeta_f)/H(K)),
\]
where the union runs over all odd primes $\ell\mid f$ and $\ell=4$.
\end{lemma}

\begin{proof}
The nontrivial element $\sigma\in\Gal(H(K)/K)\simeq\Z/2\Z$ must act
nontrivially on $H(K)$, so any lift $\hat{\sigma}$ necessarily avoids
$\Gal(\Q(\zeta_f)/H(K))$.

Because $H(K)$ is generated by adjoining square roots of suitable products of
the primes $p_i$, and since $p_3,p_4\equiv3\pmod4$, at least one such square
root (e.g.\ $\sqrt{p_1p_3}$ or $\sqrt{p_1p_4}$) lies in 
extensions whose interaction with $\Q(\zeta_\ell)$ ensures that 
$H(K)\cap\Q(\zeta_\ell)\subsetneq H(K)$ for some $\ell$.

If $H(K)\cap\Q(\zeta_\ell)\not\subset K$, then every lift of $\sigma$ avoids
$\Gal(\Q(\zeta_f)/\Q(\zeta_\ell))$.  
If, on the other hand, $H(K)\cap\Q(\zeta_\ell)\subseteq K$, one may extend
$\sigma$ to the compositum $H(K)\Q(\zeta_\ell)$ so that it restricts to
$\sigma$ on $H(K)$ and acts nontrivially on $\Q(\zeta_\ell)$.  
In either case, we obtain a lift $\hat{\sigma}$ outside the exceptional
union.
\end{proof}

\end{document}
